\chapter{Formalization and theorem correspondence}\label{app:lean}
The formal companion covers exact finite Gram algebra, disjoint block
refinement, reindexing, streamed actions, nonnegativity, coefficient order,
and arithmetic consequences of supplied enclosures. It does not establish
the Gaussian reduction, Fourier inversion, analytic remainders, spectral
post-verification algorithm, or transcendental interval-library correctness.
A theorem that assumes an enclosure does not construct that enclosure.

\begin{scopebox}
\textbf{Recorded kernel evidence.} The retained compiler record accepts
26 theorem declarations and 8 definitions, with a 34-declaration axiom audit.
The audited dependencies are confined to \texttt{propext},
\texttt{Classical.choice}, and \texttt{Quot.sound}. No \texttt{sorryAx} or
custom axiom appears. A valid arithmetic control was accepted, while false
arithmetic and an acceptance proof missing its error-budget premise were
rejected. This volume preserves that evidence and the pinned sources. It
does not claim a new Lean compilation.
\end{scopebox}

\label{app:lean-contract}
\section{How the formal model represents the engine}

The factor has a finite sample index type $I$ and an arbitrary variant index
 type $V$. A variant block is a finite set $s\subseteq V$. The factor entries
are exact real numbers, $A:I\times V\to\R$. The Lean type is
\decl{Row -> Variant -> Real}. The sample coordinates carry a
\decl{Fintype} instance. The variants carry a \decl{DecidableEq} instance to
support finite-set unions, singleton blocks, and reindexing.

The residualized weighted factor in the main text is
\[
 A=M_X G W^{1/2}.
\]
For this companion, $A$ is supplied. Proving that a particular floating-point
residualization procedure computes that mathematical factor would require
additional implementation and numerical-error reasoning. These modules do
not assume a covariance model or a distribution for the vector $v$.

There are five central finite-algebra definitions:
\begin{align*}
 \operatorname{gramEntry}(s,A,i,j)&=\sum_{k\in s}A_{ik}A_{jk},\\[6pt]
 \operatorname{matrixAction}(K,v,i)&=\sum_{j\in I}K_{ij}v_j,\\[6pt]
 \operatorname{streamedAction}(s,A,v,i)&=\sum_{k\in s}A_{ik}\sum_{j\in I}A_{jk}v_j,\\[6pt]
 \operatorname{quadraticForm}(K,v)&=\sum_{i\in I}v_i\operatorname{matrixAction}(K,v,i),\\[6pt]
 \operatorname{factorEnergy}(s,A,v)&=\sum_{k\in s}\left(\sum_{i\in I}A_{ik}v_i\right)^2.
\end{align*}
The source uses ordinary finite sums in $\R$. The factor convention is the
sample Gram $AA^\top$, matching \decl{stream_sample_gram}. It is not the dual
Gram $A^\top A$. The paper's statistic $Q=\tfrac12 v^\top Kv$ requires the
additional scalar factor $\tfrac12$ when applying these identities.

\section{Exact block refinement}

The principal refinement theorem is \decl{gramEntry_union}. Its hypotheses
are finite sets $s,t\subseteq V$ and $\operatorname{Disjoint}(s,t)$. It proves,
for all sample coordinates $i,j$,
\[
 \operatorname{gramEntry}(s\cup t,A,i,j)
 =\operatorname{gramEntry}(s,A,i,j)+\operatorname{gramEntry}(t,A,i,j).
\]
Disjointness prevents repeated columns from being counted twice. It is a
mathematical condition, not a performance assumption. The theorem applies to
empty blocks and zero factors. It does not require positive weights, since
$A$ is already supplied as a real factor.

The proof unfolds \decl{gramEntry} and invokes the finite-sum union identity.
No numerical approximation is involved. Repeated application establishes
invariance along any explicitly checked sequence of binary disjoint splits.
The source does not encode a whole scheduler or its memory allocation policy.
A scheduler correctness proof would additionally track complete coverage,
disjointness, and termination at every split.

The same refinement has formal consequences for the operator action and the
quadratic statistic:
\begin{align*}
 (K_{s\cup t}v)_i&=(K_s v)_i+(K_t v)_i,\\
 v^\top K_{s\cup t}v&=v^\top K_s v+v^\top K_t v.
\end{align*}
These are \decl{matrixAction_gram_union} and
\decl{quadraticForm_gram_union}. The independent
\decl{streamedAction_union} proves the corresponding streamed accumulation
identity directly.

The reindexing theorem \decl{gramEntry_reindex} handles a bijection
$e:V\simeq V$. It moves the block to $e(s)$ and transforms the factor to
$A'_{ik}=A_{i,e^{-1}(k)}$. The resulting Gram entry is unchanged. Permuting
only the coordinates without moving their factor columns falls outside this
statement.

\section{Streamed actions and the sum-of-squares identity}

\decl{matrixAction_gram_eq_streamed} proves
\[
 \sum_{j\in I}\left(\sum_{k\in s}A_{ik}A_{jk}\right)v_j
 =\sum_{k\in s}A_{ik}\left(\sum_{j\in I}A_{jk}v_j\right).
\]
This is the exact algebra behind multiplication by a sample Gram matrix
without constructing that matrix. Lean checks the finite-sum interchange and
associativity of real multiplication. It does not assert that a given
sequence of device operations has exact real semantics.

The next theorem, \decl{quadraticForm_gram_eq_energy}, proves
\[
 v^\top K_s v=\sum_{k\in s}\left(\sum_{i\in I}A_{ik}v_i\right)^2.
\]
It expands the action, interchanges the finite sums, factors each coordinate
sum, and identifies the product of a projection with itself. The resulting
\decl{quadraticForm_gram_nonneg} states $0\le v^\top K_s v$ for every real
vector $v$. This is positive semidefiniteness of the supplied Gram form.

These identities are useful together. Refinement preserves the represented
operator, streaming preserves its action, and the energy identity verifies
its positivity. They do not require the eigenvalues or the spectral theorem.
The full engine subsequently uses numerical rank compression and spectral
information, whose correctness contracts are separate.

\section{A complete exact algebra example}

Consider one sample coordinate and two variant columns with
\[
 A=(3\quad4),\qquad v=2.
\]
The two singleton Gram contributions are $9$ and $16$. Their disjoint union
has Gram entry $25$. Matrix multiplication gives $Kv=50$ and
$v^\top Kv=100$. The streamed action instead computes
\[
 3(3\cdot2)+4(4\cdot2)=18+32=50.
\]
The energy expression gives
\[
 (3\cdot2)^2+(4\cdot2)^2=36+64=100.
\]
Thus the engine statistic with the stated half normalization is $Q=50$.
All paths are exact real identities. This example explains the type and
normalization choices. The universal theorems establish them for arbitrary
finite sample coordinates and variant blocks.

If the two blocks overlap in one column, their sum repeats that column's
outer product. For $s=t=\{1\}$ and $A_{i1}=3$, the union has contribution $9$
while the sum of the two blocks is $18$. This is why
\decl{gramEntry_union} requires disjointness.

\section{Supplied coefficient intervals}

The separate definition
\[
 \operatorname{weightedSquares}(s,c,z)=\sum_{k\in s}c_k z_k^2
\]
models a supplied diagonal quadratic form. Its coefficients need not be
eigenvalues obtained by a proved diagonalization. It has three theorems:
\begin{itemize}
 \item \decl{weightedSquares_mono}: if $l_k\le u_k$ for every $k\in s$, then
 $\sum l_kz_k^2\le\sum u_kz_k^2$.
 \item \decl{weightedSquares_nonneg}: if $c_k\ge0$ for every $k\in s$, then
 $\sum c_kz_k^2\ge0$.
 \item \decl{weightedSquares_tail_imp}: coefficientwise order gives the
 pointwise implication $q\le\sum l_kz_k^2\Rightarrow q\le\sum u_kz_k^2$.
\end{itemize}
The first and third allow signed coefficients because squares remain
nonnegative. The second has an explicit nonnegativity hypothesis.

The pointwise implication is a prerequisite for probabilistic tail
sandwiching. It does not itself state a probability inequality or build a
coupling. To turn it into such a statement, one must supply measurable random
coordinates on a common probability space and use monotonicity of measure.
To apply it to matrix spectra, one must also establish that each true
coefficient lies inside the supplied interval.

\section{Enclosures and error ledgers}

The interval module uses two elementary predicates:
\begin{align*}
 \operatorname{Enclosed}(x,L,U)&\;:\Longleftrightarrow\;L\le x\le U,\\
 \operatorname{ErrorBound}(x,\widehat x,r)&\;:\Longleftrightarrow\;
 |x-\widehat x|\le r.
\end{align*}
The equivalence \decl{errorBound_iff_enclosed} proves
\[
 |x-\widehat x|\le r
 \quad\Longleftrightarrow\quad
 \widehat x-r\le x\le\widehat x+r.
\]
There is no extra nonnegativity hypothesis: if this predicate holds, then the
radius is necessarily nonnegative. The definition accepts arbitrary real
inputs, and the theorem states the exact logical equivalence.

For the K-COS error ledger, distinguish the ideal analytic quantity $x$, an
exact finite transform $m$, and the computed estimate $\widehat x$. If
\[
 |x-m|\le a,\qquad |m-\widehat x|\le b,
\]
then \decl{errorBound_compose} proves
\[
 |x-\widehat x|\le a+b.
\]
The proof is the triangle inequality. The hypotheses must be proved
independently. In particular, $a$ can contain a support/alias bound and an
omitted-mode bound. This theorem only composes those supplied facts.

\decl{errorBound_accept} adds the hypothesis $a+b\le\varepsilon$ and concludes
$|x-\widehat x|\le\varepsilon$. This is the arithmetic core of acceptance by a
proved budget. A sequence of numerically stable estimates, without valid
bounds, does not satisfy its hypotheses.

\section{Safe transformations of an enclosure}

The remaining interval statements are reusable exact-real rules.
\decl{enclosed_add} adds two enclosed quantities.
\decl{enclosed_scale} scales the interval by a nonnegative real factor.
\decl{enclosed_widen} permits an explicitly wider interval.
\decl{enclosed_intersect} combines two independently valid enclosures of the
same quantity using $\max$ on the lower endpoints and $\min$ on the upper
endpoints.

\decl{enclosed_positive} requires a strictly positive lower endpoint. It then
proves that the enclosed quantity is positive. This is relevant before
applying a logarithm to a corrected tilted finite sum. A lower endpoint equal
to zero does not establish positivity.

Finally, \decl{enclosed_monotone} transports an interval through a supplied
monotone function. Its hypothesis is global monotonicity on $\R$. The source
does not directly instantiate this rule with the real logarithm, which needs
a positive-domain formulation. A log-tail certificate therefore still needs
both positivity and a separate domain-correct monotonicity argument. The
interval-library implementation also requires its own correctness contract.

\section{Declaration inventory and manuscript concordance}

All names below are prefixed by \decl{KomorebiQCOS}. The inventory counts
project theorems only. Imported mathlib theorems are not counted.

{\small
\begin{longtable}{@{}p{0.49\linewidth}p{0.45\linewidth}@{}}
\toprule Declaration & Exact formal scope\\\midrule\endhead
\decl{gramEntry_empty} & Empty variant block gives zero.\\
\decl{gramEntry_singleton} & One column gives its outer product.\\
\decl{gramEntry_symm} & Real Gram entry symmetry.\\
\decl{gramEntry_union} & Disjoint binary block refinement.\\
\decl{gramEntry_reindex} & Factor-column permutation invariance.\\
\decl{matrixAction_gram_eq_streamed} & Gram action equals streamed factor action.\\
\decl{matrixAction_add} & Operator-action additivity.\\
\decl{matrixAction_gram_union} & Action refinement invariance.\\
\decl{streamedAction_union} & Streamed accumulation over disjoint blocks.\\
\decl{quadraticForm_add} & Quadratic-form additivity in the matrix.\\
\decl{quadraticForm_gram_eq_energy} & Gram form equals factor energy.\\
\decl{quadraticForm_gram_nonneg} & Positivity of the supplied Gram form.\\
\decl{quadraticForm_gram_union} & Statistic refinement invariance.\\
\decl{weightedSquares_mono} & Pointwise coefficientwise order.\\
\decl{weightedSquares_nonneg} & Positivity for nonnegative coefficients.\\
\decl{weightedSquares_tail_imp} & Pointwise tail-event implication.\\
\decl{enclosed_order} & A valid enclosure has ordered endpoints.\\
\decl{enclosed_add} & Additive interval composition.\\
\decl{enclosed_scale} & Nonnegative scaling.\\
\decl{enclosed_widen} & Valid interval widening.\\
\decl{enclosed_intersect} & Intersection of valid enclosures.\\
\decl{errorBound_iff_enclosed} & Absolute-error/enclosure equivalence.\\
\decl{errorBound_compose} & Analytic and arithmetic budget composition.\\
\decl{errorBound_accept} & Acceptance below a supplied tolerance.\\
\decl{enclosed_positive} & Positive lower endpoint proves positivity.\\
\decl{enclosed_monotone} & Supplied monotone transformation.\\
\bottomrule
\end{longtable}
}

The exact-block, hierarchical-blocking, and streamed-action sections of the
main text map to the Gram theorems under the definitions above. The
error-ledger and acceptance sections map only to the arithmetic consequences
of already valid bounds. The matrix-to-tail bridge maps to pointwise
coefficient order only. The detailed machine-readable concordance includes
source filenames and implementation anchors.

\section{Reproduction and trusted dependencies}

The companion includes the pinned toolchain and dependency manifest, two
mathematical modules, a root import, an axiom audit, controls, and retained
build logs. From the companion directory, run:
\begin{verbatim}
lake exe cache get
lake build
lake env lean AxiomAudit.lean
\end{verbatim}
Reproduction requires the pinned compiler and dependencies. Compiler
binaries and the mathlib cache are not distributed in this package.
The retained runtime record includes an executable-location adapter and its
source. It did not modify the compiler or proof kernel.

The positive and negative controls supplement the theorem build. The
axiom audit records each declaration's dependencies. Classical choice,
propositional extensionality, and quotient soundness are the reported
foundational dependencies. Any \decl{sorryAx} occurrence would indicate an
admitted proof. A clean admission scan alone does not establish acceptance.
The retained compiler logs and axiom output provide that evidence.

\section{Remaining formal obligations}

An end-to-end formal certificate would also require:
\begin{enumerate}[nosep,topsep=4pt]
 \item Gaussian whitening and central, signed, and noncentral reductions
 under their exact hypotheses.
 \item Spectral decomposition and a proved connection between the input
 matrix, its factor, and coefficient intervals.
 \item Probability monotonicity under a common measurable coupling.
 \item Exponential tilting and the transform/payoff inversion identity.
 \item Support, alias, and omitted-mode bounds with their rank and
 moment-generating-function domain conditions.
 \item Exact interval semantics for input literals, directed rounding,
 roots, complex transforms, logarithms, and transcendental operations.
 \item A verified implementation-to-bound correspondence, including
 input parsing and output endpoint conversion.
\end{enumerate}
These are separate obligations. The retained checks establish the finite
algebra and enclosure arithmetic listed in the inventory, not the
additional analytic or numerical layers.
